\haddockmoduleheading{Data.Complex}
\label{module:Data.Complex}
\haddockbeginheader
{\haddockverb\begin{verbatim}
module Data.Complex (
    Complex((:+)), realPart, imagPart, mkPolar, cis, polar, magnitude, phase,
    conjugate
  ) where\end{verbatim}}
\haddockendheader

\section{Rectangular form}
\begin{haddockdesc}
\item[\begin{tabular}{@{}l}
data Complex a
\end{tabular}]
{\haddockbegindoc
A data type representing complex numbers.\par
You can read about complex numbers \href{https://en.wikipedia.org/wiki/Complex_number}{on wikipedia}.\par
In haskell, complex numbers are represented as \haddocktt{a :+ b} which can be thought of
 as representing \(a + bi\). For a complex number \haddocktt{z}, \haddocktt{\haddockid{abs} z} is a number with the \haddockid{magnitude} of \haddocktt{z},
 but oriented in the positive real direction, whereas \haddocktt{\haddockid{signum} z}
 has the \haddockid{phase} of \haddocktt{z}, but unit \haddockid{magnitude}.
 Apart from the loss of precision due to IEEE754 floating point numbers,
 it holds that \haddocktt{z == \haddockid{abs} z * \haddockid{signum} z}.\par
Note that \haddockid{Complex}'s instances inherit the deficiencies from the type
 parameter's. For example, \haddocktt{Complex Float}'s \haddockid{Ord} instance has similar
 problems to \haddockid{Float}'s.\par
As can be seen in the examples, the \haddockid{Foldable}
 and \haddockid{Traversable} instances traverse the real part first.\par
\subsubsection*{\textbf{Examples}}
\begin{quote}
{\haddockverb\begin{verbatim}
>>> (5.0 :+ 2.5) + 6.5
11.5 :+ 2.5

\end{verbatim}}
\end{quote}
\begin{quote}
{\haddockverb\begin{verbatim}
>>> abs (1.0 :+ 1.0) - sqrt 2.0
0.0 :+ 0.0

\end{verbatim}}
\end{quote}
\begin{quote}
{\haddockverb\begin{verbatim}
>>> abs (signum (4.0 :+ 3.0))
1.0 :+ 0.0

\end{verbatim}}
\end{quote}
\begin{quote}
{\haddockverb\begin{verbatim}
>>> foldr (:) [] (1 :+ 2)
[1,2]

\end{verbatim}}
\end{quote}
\begin{quote}
{\haddockverb\begin{verbatim}
>>> mapM print (1 :+ 2)
1
2

\end{verbatim}}
\end{quote}
\enspace \emph{Constructors}\par
\haddockbeginconstrs
\haddockdecltt{=} & \haddockdecltt{!a :+ !a} & forms a complex number from its real and imaginary
 rectangular components. \\
\end{tabulary}\par}
\end{haddockdesc}
\begin{haddockdesc}
\item[\begin{tabular}{@{}l}
\end{tabular}]
\end{haddockdesc}
\begin{haddockdesc}
\item[\begin{tabular}{@{}l}
\end{tabular}]
{\haddockbegindoc
\begin{quote}
{\haddockverb\begin{verbatim}
>>> eq1 (1 :+ 2) (1 :+ 2)
True

\end{verbatim}}
\end{quote}
\begin{quote}
{\haddockverb\begin{verbatim}
>>> eq1 (1 :+ 2) (1 :+ 3)
False

\end{verbatim}}
\end{quote}}
\end{haddockdesc}
\begin{haddockdesc}
\item[\begin{tabular}{@{}l}
\end{tabular}]
{\haddockbegindoc
\begin{quote}
{\haddockverb\begin{verbatim}
>>> readPrec_to_S readPrec1 0 "(2 % 3) :+ (3 % 4)" :: [(Complex Rational, String)]
[(2 % 3 :+ 3 % 4,"")]

\end{verbatim}}
\end{quote}}
\end{haddockdesc}
\begin{haddockdesc}
\item[\begin{tabular}{@{}l}
\end{tabular}]
{\haddockbegindoc
\begin{quote}
{\haddockverb\begin{verbatim}
>>> showsPrec1 0 (2 :+ 3) ""
"2 :+ 3"

\end{verbatim}}
\end{quote}}
\end{haddockdesc}
\begin{haddockdesc}
\item[\begin{tabular}{@{}l}
instance Functor Complex\\instance Monad Complex\\instance Foldable Complex\\instance Traversable Complex\\instance Data a => Data Complex a\\instance RealFloat a => Floating Complex a\\instance Storable a => Storable Complex a\\instance RealFloat a => Num Complex a\\instance Read a => Read Complex a\\instance RealFloat a => Fractional Complex a\\instance Show a => Show Complex a\\instance Eq a => Eq Complex a\\
\end{tabular}]
\end{haddockdesc}
\begin{haddockdesc}
\item[\begin{tabular}{@{}l}
realPart :: Complex a -> a
\end{tabular}]
{\haddockbegindoc
Extracts the real part of a complex number.\par
\subsubsection*{\textbf{Examples}}
\begin{quote}
{\haddockverb\begin{verbatim}
>>> realPart (5.0 :+ 3.0)
5.0

\end{verbatim}}
\end{quote}
\begin{quote}
{\haddockverb\begin{verbatim}
>>> realPart ((5.0 :+ 3.0) * (2.0 :+ 3.0))
1.0

\end{verbatim}}
\end{quote}}
\end{haddockdesc}
\begin{haddockdesc}
\item[\begin{tabular}{@{}l}
imagPart :: Complex a -> a
\end{tabular}]
{\haddockbegindoc
Extracts the imaginary part of a complex number.\par
\subsubsection*{\textbf{Examples}}
\begin{quote}
{\haddockverb\begin{verbatim}
>>> imagPart (5.0 :+ 3.0)
3.0

\end{verbatim}}
\end{quote}
\begin{quote}
{\haddockverb\begin{verbatim}
>>> imagPart ((5.0 :+ 3.0) * (2.0 :+ 3.0))
21.0

\end{verbatim}}
\end{quote}}
\end{haddockdesc}
\section{Polar form}
\begin{haddockdesc}
\item[\begin{tabular}{@{}l}
mkPolar :: Floating a => a -> a -> Complex a
\end{tabular}]
{\haddockbegindoc
Form a complex number from \haddockid{polar} components of \haddockid{magnitude} and \haddockid{phase}.\par
\subsubsection*{\textbf{Examples}}
\begin{quote}
{\haddockverb\begin{verbatim}
>>> mkPolar 1 (pi / 4)
0.7071067811865476 :+ 0.7071067811865475

\end{verbatim}}
\end{quote}
\begin{quote}
{\haddockverb\begin{verbatim}
>>> mkPolar 1 0
1.0 :+ 0.0

\end{verbatim}}
\end{quote}}
\end{haddockdesc}
\begin{haddockdesc}
\item[\begin{tabular}{@{}l}
cis :: Floating a => a -> Complex a
\end{tabular}]
{\haddockbegindoc
\haddocktt{\haddockid{cis} t} is a complex value with \haddockid{magnitude} \haddocktt{1}
 and \haddockid{phase} \haddocktt{t} (modulo \haddocktt{2*\haddockid{pi}}).\par
\begin{quote}
{\haddockverb\begin{verbatim}
cis = mkPolar 1
\end{verbatim}}
\end{quote}
\subsubsection*{\textbf{Examples}}
\begin{quote}
{\haddockverb\begin{verbatim}
>>> cis 0
1.0 :+ 0.0

\end{verbatim}}
\end{quote}
The following examples are not perfectly zero due to \href{https://en.wikipedia.org/wiki/IEEE_754}{IEEE 754}\par
\begin{quote}
{\haddockverb\begin{verbatim}
>>> cis pi
(-1.0) :+ 1.2246467991473532e-16

\end{verbatim}}
\end{quote}
\begin{quote}
{\haddockverb\begin{verbatim}
>>> cis (4 * pi) - cis (2 * pi)
0.0 :+ (-2.4492935982947064e-16)

\end{verbatim}}
\end{quote}}
\end{haddockdesc}
\begin{haddockdesc}
\item[\begin{tabular}{@{}l}
polar :: RealFloat a => Complex a -> (a, a)
\end{tabular}]
{\haddockbegindoc
The function \haddockid{polar} takes a complex number and
 returns a (\haddockid{magnitude}, \haddockid{phase}) pair in canonical form:
 the \haddockid{magnitude} is non-negative, and the \haddockid{phase} in the range \haddocktt{(-\haddockid{pi}, \haddockid{pi}{\char 93}};
 if the \haddockid{magnitude} is zero, then so is the \haddockid{phase}.\par
\begin{quote}
{\haddockverb\begin{verbatim}
polar z = (magnitude z, phase z)\end{verbatim}}
\end{quote}
\subsubsection*{\textbf{Examples}}
\begin{quote}
{\haddockverb\begin{verbatim}
>>> polar (1.0 :+ 1.0)
(1.4142135623730951,0.7853981633974483)

\end{verbatim}}
\end{quote}
\begin{quote}
{\haddockverb\begin{verbatim}
>>> polar ((-1.0) :+ 0.0)
(1.0,3.141592653589793)

\end{verbatim}}
\end{quote}
\begin{quote}
{\haddockverb\begin{verbatim}
>>> polar (0.0 :+ 0.0)
(0.0,0.0)

\end{verbatim}}
\end{quote}}
\end{haddockdesc}
\begin{haddockdesc}
\item[\begin{tabular}{@{}l}
magnitude :: RealFloat a => Complex a -> a
\end{tabular}]
{\haddockbegindoc
The non-negative \haddockid{magnitude} of a complex number.\par
\subsubsection*{\textbf{Examples}}
\begin{quote}
{\haddockverb\begin{verbatim}
>>> magnitude (1.0 :+ 1.0)
1.4142135623730951

\end{verbatim}}
\end{quote}
\begin{quote}
{\haddockverb\begin{verbatim}
>>> magnitude (1.0 + 0.0)
1.0

\end{verbatim}}
\end{quote}
\begin{quote}
{\haddockverb\begin{verbatim}
>>> magnitude (0.0 :+ (-5.0))
5.0

\end{verbatim}}
\end{quote}}
\end{haddockdesc}
\begin{haddockdesc}
\item[\begin{tabular}{@{}l}
phase :: RealFloat a => Complex a -> a
\end{tabular}]
{\haddockbegindoc
The \haddockid{phase} of a complex number, in the range \haddocktt{(-\haddockid{pi}, \haddockid{pi}{\char 93}}.
 If the \haddockid{magnitude} is zero, then so is the \haddockid{phase}.\par
\subsubsection*{\textbf{Examples}}
\begin{quote}
{\haddockverb\begin{verbatim}
>>> phase (0.5 :+ 0.5) / pi
0.25

\end{verbatim}}
\end{quote}
\begin{quote}
{\haddockverb\begin{verbatim}
>>> phase (0 :+ 4) / pi
0.5

\end{verbatim}}
\end{quote}}
\end{haddockdesc}
\section{Conjugate}
\begin{haddockdesc}
\item[\begin{tabular}{@{}l}
conjugate :: Num a => Complex a -> Complex a
\end{tabular}]
{\haddockbegindoc
The \haddockid{conjugate} of a complex number.\par
\begin{quote}
{\haddockverb\begin{verbatim}
conjugate (conjugate x) = x\end{verbatim}}
\end{quote}
\subsubsection*{\textbf{Examples}}
\begin{quote}
{\haddockverb\begin{verbatim}
>>> conjugate (3.0 :+ 3.0)
3.0 :+ (-3.0)

\end{verbatim}}
\end{quote}
\begin{quote}
{\haddockverb\begin{verbatim}
>>> conjugate ((3.0 :+ 3.0) * (2.0 :+ 2.0))
0.0 :+ (-12.0)

\end{verbatim}}
\end{quote}}
\end{haddockdesc}
\section{Specification}
\begin{quote}
{\haddockverb\begin{verbatim}
module Data.Complex(Complex((:+)), realPart, imagPart, conjugate, mkPolar,
                    cis, polar, magnitude, phase)  where

infix  6  :+

data  (RealFloat a)     => Complex a = !a :+ !a  deriving (Eq,Read,Show)


realPart, imagPart :: (RealFloat a) => Complex a -> a
realPart (x:+y)        =  x
imagPart (x:+y)        =  y

conjugate      :: (RealFloat a) => Complex a -> Complex a
conjugate (x:+y) =  x :+ (-y)

mkPolar                :: (RealFloat a) => a -> a -> Complex a
mkPolar r theta        =  r ⋆ cos theta :+ r ⋆ sin theta

cis            :: (RealFloat a) => a -> Complex a
cis theta      =  cos theta :+ sin theta

polar          :: (RealFloat a) => Complex a -> (a,a)
polar z                =  (magnitude z, phase z)

magnitude :: (RealFloat a) => Complex a -> a
magnitude (x:+y) =  scaleFloat k
                   (sqrt ((scaleFloat mk x)^2 + (scaleFloat mk y)^2))
                  where k  = max (exponent x) (exponent y)
                        mk = - k

phase :: (RealFloat a) => Complex a -> a
phase (0 :+ 0) = 0
phase (x :+ y) = atan2 y x


instance  (RealFloat a) => Num (Complex a)  where
    (x:+y) + (x':+y') =  (x+x') :+ (y+y')
    (x:+y) - (x':+y') =  (x-x') :+ (y-y')
    (x:+y) ⋆ (x':+y') =  (x⋆x'-y⋆y') :+ (x⋆y'+y⋆x')
    negate (x:+y)     =  negate x :+ negate y
    abs z             =  magnitude z :+ 0
    signum 0          =  0
    signum z(x:+y)   =  x\emph{r :+ y}r  where r = magnitude z
    fromInteger n     =  fromInteger n :+ 0

instance  (RealFloat a) => Fractional (Complex a)  where
    (x:+y) \emph{ (x':+y') =  (x⋆x''+y⋆y'') } d :+ (y⋆x''-x⋆y'') / d
                         where x'' = scaleFloat k x'
                               y'' = scaleFloat k y'
                               k   = - max (exponent x') (exponent y')
                               d   = x'⋆x'' + y'⋆y''

    fromRational a    =  fromRational a :+ 0

instance  (RealFloat a) => Floating (Complex a)       where
    pi             =  pi :+ 0
    exp (x:+y)     =  expx ⋆ cos y :+ expx ⋆ sin y
                      where expx = exp x
    log z          =  log (magnitude z) :+ phase z

    sqrt 0         =  0
    sqrt z(x:+y)  =  u :+ (if y < 0 then -v else v)
                      where (u,v) = if x < 0 then (v',u') else (u',v')
                            v'    = abs y / (u'⋆2)
                            u'    = sqrt ((magnitude z + abs x) / 2)

    sin (x:+y)     =  sin x ⋆ cosh y :+ cos x ⋆ sinh y
    cos (x:+y)     =  cos x ⋆ cosh y :+ (- sin x ⋆ sinh y)
    tan (x:+y)     =  (sinx⋆coshy:+cosx⋆sinhy)/(cosx⋆coshy:+(-sinx⋆sinhy))
                      where sinx  = sin x
                            cosx  = cos x
                            sinhy = sinh y
                            coshy = cosh y

    sinh (x:+y)    =  cos y ⋆ sinh x :+ sin  y ⋆ cosh x
    cosh (x:+y)    =  cos y ⋆ cosh x :+ sin y ⋆ sinh x
    tanh (x:+y)    =  (cosy⋆sinhx:+siny⋆coshx)/(cosy⋆coshx:+siny⋆sinhx)
                      where siny  = sin y
                            cosy  = cos y
                            sinhx = sinh x
                            coshx = cosh x

    asin z(x:+y)  =  y':+(-x')
                      where  (x':+y') = log (((-y):+x) + sqrt (1 - z⋆z))
    acos z(x:+y)  =  y'':+(-x'')
                      where (x'':+y'') = log (z + ((-y'):+x'))
                            (x':+y')   = sqrt (1 - z⋆z)
    atan z@(x:+y)  =  y':+(-x')
                      where (x':+y') = log (((1-y):+x) / sqrt (1+z⋆z))

    asinh z        =  log (z + sqrt (1+z⋆z))
    acosh z        =  log (z + (z+1) ⋆ sqrt ((z-1)/(z+1)))
    atanh z        =  log ((1+z) / sqrt (1-z⋆z))
\end{verbatim}}
\end{quote}